\section{Определение евклидовых и унитарных пространств. Примеры. Норма вектора. Неравенство Коши-Буняковского. Неравенство треугольника}

\begin{definition}
Линейное пространство $V$ над полем $\mathbb{R}$ называется \textbf{евклидовым}, если каждой паре векторов $x, y \in V$ поставлено в соответствие действительное число $(x, y)$, называемое \textbf{скалярным произведением}, удовлетворяющее следующим аксиомам:
\begin{enumerate}
    \item $(x, y) = (y, x)$ (симметричность);
    \item $(x_1 + x_2, y) = (x_1, y) + (x_2, y)$ (аддитивность по первому аргументу);
    \item $(\lambda x, y) = \lambda (x, y)$ для любого $\lambda \in \mathbb{R}$ (однородность по первому аргументу);
    \item $(x, x) > 0$ при $x \neq 0$, и $(x, x) = 0 \iff x = 0$ (положительная определенность).
\end{enumerate}
\end{definition}

\begin{definition}
Линейное пространство $V$ над полем $\mathbb{C}$ называется \textbf{унитарным}, если скалярное произведение $(x, y) \in \mathbb{C}$ удовлетворяет аксиомам:
\begin{enumerate}
    \item $(x, y) = \overline{(y, x)}$ (эрмитова симметричность, где черта обозначает комплексное сопряжение);
    \item $(x_1 + x_2, y) = (x_1, y) + (x_2, y)$;
    \item $(\lambda x, y) = \lambda (x, y)$ для любого $\lambda \in \mathbb{C}$;
    \item $(x, x) \in \mathbb{R}$, причем $(x, x) > 0$ при $x \neq 0$, и $(x, x) = 0 \iff x = 0$.
\end{enumerate}
\end{definition}

\begin{example}[Примеры]
\begin{enumerate}
    \item Пространство $\mathbb{R}^n$ со скалярным произведением $(x, y) = \sum_{i=1}^n x_i y_i$.
    \item Пространство $\mathbb{C}^n$ со скалярным произведением $(x, y) = \sum_{i=1}^n x_i \overline{y_i}$.
    \item Пространство непрерывных функций $C([a, b])$ со скалярным произведением $(f, g) = \int_a^b f(t) g(t) \, dt$ (евклидово) или $\int_a^b f(t) \overline{g(t)} \, dt$ (унитарное).
\end{enumerate}
\end{example}

\begin{definition}
\textbf{Нормой} (или длиной) вектора $x$ в евклидовом или унитарном пространстве называется неотрицательное число:
\[
\|x\| = \sqrt{(x, x)}.
\]
\end{definition}

\begin{theorem}[Неравенство Коши-Буняковского]
Для любых векторов $x, y$ в евклидовом или унитарном пространстве справедливо неравенство:
\[
|(x, y)| \le \|x\| \cdot \|y\|.
\]
\end{theorem}

\begin{proof}
Если $y = 0$, то $(x, 0) = 0$ и $\|y\| = 0$, неравенство превращается в $0 \le 0$, что верно.
Пусть $y \neq 0$. Рассмотрим вектор $z = x - \lambda y$, где $\lambda$ --- некоторое число. В силу положительной определенности скалярного произведения:
\[
(z, z) = (x - \lambda y, x - \lambda y) \ge 0.
\]
Раскроем скобки, используя свойства скалярного произведения. Для унитарного пространства:
\[
(x, x) - \overline{\lambda}(x, y) - \lambda(y, x) + |\lambda|^2 (y, y) \ge 0.
\]
Выберем $\lambda = \frac{(x, y)}{(y, y)}$ (это возможно, так как $(y, y) = \|y\|^2 > 0$). Тогда $\overline{\lambda} = \frac{\overline{(x, y)}}{(y, y)}$. Подставим это значение:
\[
\|x\|^2 - \frac{\overline{(x, y)}(x, y)}{\|y\|^2} - \frac{(x, y)\overline{(x, y)}}{\|y\|^2} + \frac{|(x, y)|^2}{\|y\|^4} \|y\|^2 \ge 0.
\]
Упрощая, получаем:
\[
\|x\|^2 - \frac{|(x, y)|^2}{\|y\|^2} - \frac{|(x, y)|^2}{\|y\|^2} + \frac{|(x, y)|^2}{\|y\|^2} \ge 0 \quad \Rightarrow \quad \|x\|^2 - \frac{|(x, y)|^2}{\|y\|^2} \ge 0.
\]
Умножая на $\|y\|^2 > 0$, получаем $|(x, y)|^2 \le \|x\|^2 \|y\|^2$. Извлекая квадратный корень, приходим к требуемому неравенству.
\hfill$\blacksquare$
\end{proof}

\begin{theorem}[Неравенство треугольника]
Для любых векторов $x, y$ в евклидовом или унитарном пространстве справедливо неравенство:
\[
\|x + y\| \le \|x\| + \|y\|.
\]
\end{theorem}

\begin{proof}
Рассмотрим квадрат нормы суммы векторов:
\[
\|x + y\|^2 = (x + y, x + y) = (x, x) + (x, y) + (y, x) + (y, y).
\]
В евклидовом пространстве $(x, y) = (y, x)$, поэтому $(x, y) + (y, x) = 2(x, y) \le 2|(x, y)|$.
В унитарном пространстве $(y, x) = \overline{(x, y)}$, поэтому $(x, y) + (y, x) = 2 \operatorname{Re}(x, y) \le 2|(x, y)|$.
В обоих случаях:
\[
\|x + y\|^2 \le \|x\|^2 + 2|(x, y)| + \|y\|^2.
\]
Применяя неравенство Коши-Буняковского $|(x, y)| \le \|x\| \|y\|$, получаем:
\[
\|x + y\|^2 \le \|x\|^2 + 2\|x\|\|y\| + \|y\|^2 = (\|x\| + \|y\|)^2.
\]
Поскольку нормы неотрицательны, извлекая квадратный корень из обеих частей, получаем $\|x + y\| \le \|x\| + \|y\|$.
\hfill$\blacksquare$
\end{proof}
